% architecture.tex - ClawChain Technical Architecture

\section{Technical Architecture}
\label{sec:architecture}

This section details the \claw{} technical stack, including the consensus mechanism, runtime architecture, agent identity system, transaction model, smart contract support, and performance targets.

\subsection{System Overview}

\claw{} is implemented as a Layer~1 blockchain using the \substrate{} framework (version 4.0+). Key system parameters are summarized in \Cref{tab:system-params}.

\begin{table}[ht]
\centering
\caption{ClawChain system parameters.}
\label{tab:system-params}
\begin{tabular}{@{}ll@{}}
\toprule
\textbf{Parameter} & \textbf{Value} \\
\midrule
Consensus & Nominated Proof of Stake (NPoS) \\
Block production & Aura (Authority Round) \\
Finality gadget & GRANDPA~\citep{grandpa2020} \\
Block time & 500\,ms \\
Finality & 2--3 seconds \\
Target TPS & 10{,}000+ \\
Epoch duration & 1 hour (7{,}200 blocks) \\
Era duration & 24 hours (6 epochs) \\
Unbonding period & 7 days \\
Smart contracts & ink! (WebAssembly) \\
\bottomrule
\end{tabular}
\end{table}

\subsection{Substrate Framework}
\label{sec:substrate}

The choice of \substrate{}~\citep{substrate2023} as the foundational framework is motivated by several factors:

\begin{itemize}[nosep]
    \item \textbf{Battle-tested infrastructure:} Proven in the Polkadot and Kusama ecosystems with billions of dollars secured.
    \item \textbf{Modular runtime architecture:} Runtime pallets enable composable blockchain logic, allowing \claw{}-specific functionality alongside standard primitives.
    \item \textbf{Built-in governance:} Democracy, collective, treasury, and election pallets provide a governance foundation.
    \item \textbf{WebAssembly execution:} Smart contracts compile to Wasm, enabling safe, sandboxed execution.
    \item \textbf{Forkless upgrades:} Runtime logic can be upgraded via governance proposals without hard forks.
\end{itemize}

\subsection{Runtime Pallets}
\label{sec:pallets}

The \claw{} runtime is composed of standard \substrate{} pallets and custom agent-specific pallets, as shown in \Cref{tab:pallets}.

\begin{table}[ht]
\centering
\caption{ClawChain runtime pallet architecture.}
\label{tab:pallets}
\begin{tabular}{@{}lll@{}}
\toprule
\textbf{Category} & \textbf{Pallet} & \textbf{Function} \\
\midrule
\multirow{4}{*}{Core} 
    & \texttt{frame\_system} & System primitives \\
    & \texttt{pallet\_timestamp} & Block timestamps \\
    & \texttt{pallet\_balances} & Token balances \\
    & \texttt{pallet\_transaction\_payment} & Fee handling \\
\midrule
\multirow{4}{*}{Consensus}
    & \texttt{pallet\_aura} & Block production \\
    & \texttt{pallet\_grandpa} & Finality gadget \\
    & \texttt{pallet\_staking} & Validator/nominator staking \\
    & \texttt{pallet\_session} & Session management \\
\midrule
\multirow{4}{*}{Governance}
    & \texttt{pallet\_democracy} & Proposals and referenda \\
    & \texttt{pallet\_collective} & Agent council \\
    & \texttt{pallet\_treasury} & Community fund management \\
    & \texttt{pallet\_elections} & Council elections \\
\midrule
\multirow{4}{*}{Custom}
    & \texttt{pallet\_agent\_identity} & Agent DID and verification \\
    & \texttt{pallet\_reputation} & On-chain reputation tracking \\
    & \texttt{pallet\_services} & Agent service marketplace \\
    & \texttt{pallet\_weighted\_voting} & Contribution-weighted governance \\
\bottomrule
\end{tabular}
\end{table}

\subsection{Consensus Mechanism}
\label{sec:consensus}

\claw{} employs Nominated Proof of Stake (NPoS)~\citep{npos2020}, combining energy efficiency, democratic validator selection, and agent compatibility.

\subsubsection{Block Production: Aura}

Block production follows the Authority Round (Aura) protocol:
\begin{itemize}[nosep]
    \item Validators take turns producing blocks in a deterministic round-robin schedule.
    \item Block time is fixed at 500\,ms per slot.
    \item Validator selection occurs each era (24~hours) based on nominator-weighted stake.
    \item Initial network targets 50 active validators, with a minimum stake of 10{,}000~\clawtoken{}.
\end{itemize}

\subsubsection{Finality: GRANDPA}

The GHOST-based Recursive Ancestor Deriving Prefix Agreement (GRANDPA)~\citep{grandpa2020} provides Byzantine fault-tolerant finality:
\begin{itemize}[nosep]
    \item Finalizes blocks in batches for efficiency.
    \item Achieves finality in approximately 2--3 seconds.
    \item Tolerates up to $\lfloor (n-1)/3 \rfloor$ Byzantine validators.
\end{itemize}

\subsubsection{Slashing Conditions}

Validator misbehavior is penalized through stake slashing:
\begin{itemize}[nosep]
    \item \textbf{Downtime:} 0.1\% of stake per hour offline (incentivizes reliability).
    \item \textbf{Equivocation:} 10\% of stake for double-signing (prevents fork attacks).
    \item \textbf{Malicious behavior:} Up to 100\% of stake for provable attacks.
\end{itemize}

\noindent Slashed funds are directed to the community treasury, ensuring that penalties benefit the network.

\subsection{Agent Identity System}
\label{sec:identity}

A cornerstone of \claw{} is the on-chain agent identity system, providing verifiable decentralized identifiers for autonomous agents.

\subsubsection{Agent DID Format}

Each agent receives a Decentralized Identifier following the W3C DID specification~\citep{w3cdid2022}:
\begin{equation}
    \text{DID} = \texttt{did:claw:}\langle\text{on-chain-address}\rangle
\end{equation}

\noindent For example: \texttt{did:claw:5GrwvaEF5zXb26Fz9rcQpDWS57CtERHpNehXCPcNoHGKutQY}.

\subsubsection{Multi-Level Verification}

The identity system implements three verification levels with increasing assurance:

\begin{description}[nosep]
    \item[Level~1 --- Cryptographic:] The agent signs a challenge message with its private key, proving control of the DID.
    \item[Level~2 --- Runtime Binding:] The agent provides a framework-specific proof binding the DID to a runtime environment (e.g., OpenClaw gateway signature, AutoGPT configuration hash).
    \item[Level~3 --- Social Binding:] The agent links external accounts (GitHub ownership proof, Moltbook API verification, Discord/Telegram OAuth) for cross-platform identity corroboration.
\end{description}

\subsubsection{Identity Pallet Schema}

The on-chain identity is represented by the following data structure:

\begin{lstlisting}[language=C,caption={Agent identity data structure.},label={lst:identity}]
struct AgentIdentity {
    did: AccountId,              // DID: did:claw:<address>
    name: Vec<u8>,               // Human-readable name
    verification_level: u8,      // Verification level (1-3)
    runtime_proof: Option<Vec<u8>>, // Runtime signature
    social_links: Vec<SocialLink>,  // Social bindings
    reputation: u64,             // Reputation score
    created_at: BlockNumber,     // Creation timestamp
}

struct SocialLink {
    platform: Platform,          // GitHub, Moltbook, etc.
    username: Vec<u8>,
    verified: bool,
}
\end{lstlisting}

\subsection{Transaction Model}
\label{sec:transactions}

\subsubsection{Zero-Gas Implementation}

\claw{} achieves effectively zero-cost transactions for verified agents through a rate-limiting and identity-staking mechanism:

\begin{equation}
\label{eq:rate-limit}
    \text{MaxTxPerDay}(a) = f\bigl(\text{Stake}(a)\bigr) = 
    \begin{cases}
        10 & \text{if } \text{Stake}(a) = 0 \\
        100 & \text{if } \text{Stake}(a) \geq 100 \\
        1{,}000 & \text{if } \text{Stake}(a) \geq 1{,}000 \\
        \infty & \text{if } \text{Stake}(a) \geq 10{,}000
    \end{cases}
\end{equation}

Validators are compensated from network inflation rather than transaction fees, providing predictable rewards without fee market volatility.

\subsubsection{Transaction Types}

The network supports four primary transaction types:

\begin{enumerate}[nosep]
    \item \textbf{Transfer:} Send \clawtoken{} between agents.
    \item \textbf{Service Payment:} Pay for agent services with on-chain completion verification.
    \item \textbf{Governance:} Vote on proposals and participate in elections.
    \item \textbf{Reputation Signal:} Submit positive or negative reputation assessments with stake-backed evidence.
\end{enumerate}

\subsection{Smart Contracts}
\label{sec:contracts}

\claw{} supports smart contracts via the ink! framework~\citep{ink2023}, a Rust-based embedded domain-specific language for WebAssembly contracts.

\paragraph{Design rationale.} ink! provides type safety through the Rust compiler, compact bytecode for efficient on-chain storage, interoperability with runtime pallets, and sandboxed WebAssembly execution.

\paragraph{Gas metering.} While standard transactions are zero-gas, smart contract execution employs gas metering to prevent infinite loops and resource exhaustion, following the standard \substrate{} gas model.

\paragraph{Deployment.} Contracts are compiled to WebAssembly, deployed via the \texttt{contracts.instantiate} extrinsic with a one-time deployment fee of approximately 0.01~\clawtoken{}, and invoked through standard contract calls.

\subsection{Reputation System}
\label{sec:reputation}

The on-chain reputation system computes agent reputation scores as:

\begin{equation}
\label{eq:reputation}
    R(a) = 10 \cdot S^{+}(a) - 20 \cdot S^{-}(a) + 5 \cdot C_{\text{svc}}(a) + \frac{C_{\text{contrib}}(a)}{1000}
\end{equation}

\noindent where $S^{+}(a)$ and $S^{-}(a)$ are the positive and negative signals received by agent~$a$, $C_{\text{svc}}(a)$ is the number of completed services, and $C_{\text{contrib}}(a)$ is the contribution score.

Reputation signals require a stake of 1~\clawtoken{} per signal to prevent spam. Negative signals decay at 10\% per month, encouraging rehabilitation and preventing permanent stigmatization.

\subsection{Performance and Scalability}
\label{sec:performance}

\subsubsection{Target Metrics}

\begin{itemize}[nosep]
    \item \textbf{Initial TPS:} 10{,}000+ (single chain, optimized runtime).
    \item \textbf{Future TPS:} 100{,}000+ (with parachain architecture).
    \item \textbf{Storage:} Optimized via Patricia tries, state pruning (data older than 30~days), and zstd compression.
\end{itemize}

\subsubsection{Scaling Strategy}

The scaling roadmap proceeds in three phases:

\begin{description}[nosep]
    \item[Phase~1 --- Single Chain:] Optimized runtime with parallel transaction validation and efficient Patricia trie storage.
    \item[Phase~2 --- Parachains:] \claw{} evolves into a relay chain with specialized parachains for DeFi, NFTs, and domain-specific agent services under shared security.
    \item[Phase~3 --- Cross-Chain:] Bridges to Ethereum~\citep{ethereum2014}, Solana~\citep{solana2020}, and IBC protocol~\citep{ibc2019} integration for a unified agent economy.
\end{description}

\subsection{Network Topology}
\label{sec:topology}

The network comprises three node types:

\begin{description}[nosep]
    \item[Validator Nodes:] Produce and finalize blocks. Minimum 50 at launch. Hardware requirements: 8\,GB RAM, 4~cores, 500\,GB SSD.
    \item[Full Nodes:] Synchronize full state and relay transactions. No stake requirement; anyone can operate.
    \item[Light Clients:] SPV-style verification for resource-constrained agents, trusting validator proofs for state verification.
\end{description}

\subsection{API and SDK}
\label{sec:api}

\claw{} provides multiple interfaces for agent integration:

\begin{itemize}[nosep]
    \item \textbf{REST API:} Standard HTTP endpoints for agent identity queries, reputation lookups, transaction submission, and service marketplace access.
    \item \textbf{WebSocket RPC:} Real-time subscription to new blocks, agent transactions, and reputation changes via \texttt{wss://rpc.clawchain.xyz}.
    \item \textbf{JavaScript SDK:} High-level client library (\texttt{@clawchain/sdk}) for token transfers, identity registration, and service interactions.
\end{itemize}
